% Part: model-theory
% Chapter: models-of-arithmetic
% Section: introduction

\documentclass[../../../include/open-logic-section]{subfiles}

\begin{document}

\olfileid[pt]{mod}{mar}{int}
\section{Introdução}

O \emph{modelo padrão} da aritmética é a
!!{structure}~$\Struct{N}$ com $\Domain{N} = \Nat$, na qual
$\Obj{0}$, $\prime$, $+$, $\times$ e $<$ são interpretados como se
espera. Isto é, $\Obj{0}$ é $0$, $\prime$ é a função sucessor,
$+$ é interpretado como a adição e $\times$ como a multiplicação
dos números em~$\Nat$. Mais precisamente,
\begin{align*}
  \Assign{\Obj{0}}{N} & = 0\\
  \Assign{\prime}{N}(n) & = n + 1\\
  \Assign{+}{N}(n, m) & = n + m\\
  \Assign{\times}{N}(n, m) & = nm
\end{align*}
É claro que existem estruturas para $\Lang{L_A}$ cujos domínios são
diferentes de~$\Nat$. Por exemplo, podemos tomar $\Struct{M}$ com domínio
$\Domain{M} = \{a\}^*$ (as sequências finitas do único
símbolo~$a$, isto é, $\emptyset$, $a$, $aa$, $aaa$, \dots) e as
interpretações
\begin{align*}
  \Assign{\Obj{0}}{M} & = \emptyset\\
  \Assign{\prime}{M}(s) & = s \concat a\\
  \Assign{+}{M}(n, m) & = a^{n + m}\\
  \Assign{\times}{M}(n, m) & = a^{nm}
\end{align*}
Essas duas estruturas são ``essencialmente iguais'', no sentido de que
a única diferença está nos !!{element}s dos !!{domain}s, e não na maneira
como os !!{element}s dos !!{domain}s se relacionam entre si por meio
das funções de interpretação. Dizemos que as duas !!{structure}s são
\emph{isomorfas}.

É uma consequência simples do teorema da compacidade que toda teoria
verdadeira em~$\Struct{N}$ também tem modelos que não são isomorfos
a~$\Struct{N}$.  Tais estruturas são chamadas \emph{não padrão}.  O
interessante é que, enquanto os !!{element}s de um modelo padrão (isto
é, $\Struct{N}$, mas também todas as !!{structure}s isomorfas a ele) se
esgotam nos valores dos numerais padrão~$\num{n}$, isto é,
\[
\Domain{N} = \Setabs{\Value{\num{n}}{N}}{n \in \Nat}
\]
isso não ocorre nos modelos não padrão: se $\Struct{M}$ é não padrão,
então existe pelo menos um $x \in \Domain{M}$ tal que
$x \neq \Value{\num{n}}{M}$ para todo~$n$.

Esses elementos não padrão são bastante interessantes: são ``números
naturais infinitos.'' Mas sua existência também explica, em certo
sentido, os fenômenos de incompletude.  Considere, por exemplo, a
sentença de consistência da aritmética de Peano, $\OCon[\Th{PA}]$, isto
é, $\lnot \lexists[x][\OPrf[\Th{PA}](x, \gn{\lfalse})]$. Como $\Th{PA}$
não demonstra $\OCon[\Th{PA}]$ nem $\lnot \OCon[\Th{PA}]$, qualquer uma
das duas pode ser adicionada consistentemente a $\Th{PA}$. Como
$\Th{PA}$ é consistente, $\Sat{N}{\OCon[\Th{PA}]}$ e, consequentemente,
$\Sat/{N}{\lnot \OCon[\Th{PA}]}$.  Portanto, $\Struct{N}$ \emph{não} é
um modelo de $\Th{PA} \cup \{\lnot \OCon[\Th{PA}]\}$, e todos os seus
modelos devem ser não padrão. Modelos de $\Th{PA} \cup \{\lnot
\OCon[\Th{PA}]\}$ devem conter algum !!{element} que sirva como
testemunha e torne $\lexists[x][\OPrf[\Th{PA}](\gn{\lfalse})]$
verdadeira, isto é, um número de G\"odel de !!a{derivation} de uma
contradição a partir de~$\Th{PA}$.  Tal !!a{element} não pode ser
padrão---pois $\Th{PA} \Proves \lnot
\OPrf[\Th{PA}](\num{n}, \gn{\lfalse})$ para todo~$n$.

\end{document}
